\appendix

%==========================================================================
\section{The three-term identity}\label{app:identity}
%==========================================================================
Let $d_k=(h_k+h_{k+1})/2$ be the distance between neighbouring midpoints. For every quadratic spline on the
grid, with nodal values $U_k=s(m_k)$ and cell values $c_k=s''|_{I_k}$, the interior rows of~\eqref{eq:identity}
read
\begin{equation}\label{eq:interior}
\frac{U_{k+1}-U_k}{d_k}-\frac{U_k-U_{k-1}}{d_{k-1}}
=\frac{h_{k-1}^2}{8d_{k-1}}c_{k-1}+\Big(h_k-\frac{h_k^2}{8d_{k-1}}-\frac{h_k^2}{8d_k}\Big)c_k+\frac{h_{k+1}^2}{8d_k}c_{k+1}.
\end{equation}
The left-hand side is the standard difference of slopes between midpoints, and the right-hand side a
weighted average of $u''$ over three cells. At the left end, with a prescribed value $g_0=s(a)$ or a zero
slope,
\begin{align}
\frac{U_2-U_1}{d_1}-\frac{2(U_1-g_0)}{h_1}&=\Big(\frac{3h_1}{4}-\frac{h_1^2}{8d_1}\Big)c_1+\frac{h_2^2}{8d_1}c_2 &&\text{(value prescribed)},\label{eq:dirrow}\\
\frac{U_2-U_1}{d_1}&=\Big(h_1-\frac{h_1^2}{8d_1}\Big)c_1+\frac{h_2^2}{8d_1}c_2 &&\text{(zero slope)},\label{eq:neurow}
\end{align}
and symmetrically at the right end. The matrix $\mathbf A$ has positive entries and is strictly diagonally dominant
by columns, so it is always invertible. The collocation problem $s''(m_k)=f(m_k)$ has a unique solution if a
value is prescribed at one or both ends; if the slope vanishes at both ends, it is solvable if and only if
$\sum_kh_kf(m_k)=0$, and unique up to a constant (Result~1 of the Supplementary Material).

Slopes follow by summing the cell values of $u''$. With the left-end slope
$\sigma_0=2(U_1-g_0)/h_1-c_1h_1/4$ (zero for a zero-slope end),
\begin{equation}\label{eq:D1}
s'(m_k)=\sigma_0+\sum_{j<k}h_jc_j+\tfrac12h_kc_k,\qquad s'(b)=\sigma_0+\sum_kh_kc_k,
\end{equation}
which, with $\mathbf c=\mathbf D_2\mathbf U+\mathbf e$, gives the first-derivative matrix $\mathbf D_1$ and the end slopes of Algorithm~1.

%==========================================================================
\section{Time stepping for the phase-field equations}\label{app:sav}
%==========================================================================
Let $\mathbf u^n$ be the vector of nodal values at time step $n$, $\mathbf K$ the discrete $-\Delta$, $\mathbf G=\mathbf I$ (Allen--Cahn) or $\mathbf G=\mathbf K$ (Cahn--Hilliard), $\bar{\mathbf u}=2\mathbf u^n-\mathbf u^{n-1}$, and
$r$ a scalar auxiliary variable~\cite{shen2018scalar}. The stabilized SAV-BDF2 scheme reads
\begin{align}
\frac{3\mathbf u^{n+1}-4\mathbf u^n+\mathbf u^{n-1}}{2\Delta t}&=-\mathbf G\boldsymbol\mu^{n+1},\label{eq:sav1}\\
\boldsymbol\mu^{n+1}&=\varepsilon^2\mathbf K\mathbf u^{n+1}+S(\mathbf u^{n+1}-\bar{\mathbf u})+r^{n+1}\bar{\mathbf b},\label{eq:sav1b}\\
3r^{n+1}-4r^n+r^{n-1}&=\tfrac12\langle\bar{\mathbf b},3\mathbf u^{n+1}-4\mathbf u^n+\mathbf u^{n-1}\rangle_h,\label{eq:sav2}
\end{align}
where $\bar{\mathbf b}=f(\bar{\mathbf u})/\sqrt{E_1(\bar{\mathbf u})+C}$ is the scaled nonlinearity, $f(u)=u^3-u$ applied entrywise,
$E_1(\mathbf u)=\langle\frac14(\mathbf u^2-1)^2,1\rangle_h$, $C>0$, and $S\ge0$ a stabilization constant. On a uniform
grid, for every time step and every $S\ge0$, the scheme has a unique solution, the Cahn--Hilliard scheme
conserves the discrete mass exactly, and a modified energy never increases (Result~5 of the Supplementary
Material). Algorithm~2 lists one step.

\begin{algorithm}[t]\small\hrule\smallskip
\noindent\textbf{Algorithm 2: one SAV-BDF2 time step} (\texttt{ihaar/phasefield.py}).
Precomputed: the transform $\mathcal T$ (DCT-II, or the eigenvectors $\mathbf V$), the symbols $\hat{\boldsymbol\lambda}$ of
$-\mathbf K$, $\hat{\mathbf G}=\mathbf 1$ (AC) or $\hat{\boldsymbol\lambda}$ (CH), and $\hat{\mathbf S}_2=3+2\Delta t\,\hat{\mathbf G}(\varepsilon^2\hat{\boldsymbol\lambda}+S)$ (entrywise).
Stored from previous steps: $\hat{\mathbf u}^n$, $\hat{\mathbf u}^{n-1}$, $\mathbf u^n$, $\mathbf u^{n-1}$, $r^n$, $r^{n-1}$.
\step{1} Compute $\bar{\mathbf u}=2\mathbf u^n-\mathbf u^{n-1}$ and $\bar{\mathbf b}=f(\bar{\mathbf u})/\sqrt{E_1(\bar{\mathbf u})+C}$, and transform, $\hat{\bar{\mathbf b}}=\mathcal T\bar{\mathbf b}$ (one forward transform).
\step{2} Compute $\hat{\mathbf g}=4\hat{\mathbf u}^n-\hat{\mathbf u}^{n-1}$, $\hat{\mathbf p}=(\hat{\mathbf g}+2\Delta t\,S\hat{\mathbf G}\,\hat{\bar{\mathbf u}})/\hat{\mathbf S}_2$ and $\hat{\mathbf q}=-2\Delta t\,\hat{\mathbf G}\hat{\bar{\mathbf b}}/\hat{\mathbf S}_2$ (entrywise).
\step{3} Update the scalar, $r^{n+1}=\big(4r^n-r^{n-1}+\tfrac12\langle\hat{\bar{\mathbf b}},3\hat{\mathbf p}-\hat{\mathbf g}\rangle\big)/\big(3-\tfrac32\langle\hat{\bar{\mathbf b}},\hat{\mathbf q}\rangle\big)$, with inner products evaluated in the orthonormal transform domain.
\step{4} Set $\hat{\mathbf u}^{n+1}=\hat{\mathbf p}+r^{n+1}\hat{\mathbf q}$ and $\mathbf u^{n+1}=\mathcal T^{-1}\hat{\mathbf u}^{n+1}$ (one inverse transform).
\smallskip\hrule
\end{algorithm}

